\section{Credit Assignment and Optimization with Verbal Signals}
\label{sec:credit}

Credit assignment asks which parts of an agent's behavior deserve blame or praise for a delayed outcome. Language makes this problem harder because an LLM action may span thousands of tokens and several turns, yet it also offers a distinctive solution: a critique can name the faulty step and the needed repair, information that no scalar carries. We organize the literature by the granularity at which credit lands (outcome, step, turn, token, or memory operation) and by its carrier (redistributed scalar, model-internal attribution, or critique text). These choices determine which optimizer can consume the signal and which assumptions must hold.

\begin{table}[t]
\centering
\caption{Families of credit-assignment mechanisms in VRL, ordered by how much linguistic content the credit carrier consumes.}
\label{tab:credit-families}
\small
\begin{tabular}{@{}p{0.19\textwidth}p{0.30\textwidth}p{0.13\textwidth}p{0.26\textwidth}@{}}
\toprule
Family & Credit carrier & Granularity & Representative work \\
\midrule
Temporal abstraction & scalar reward over macro actions or turns & span, turn & \citep{chai2024marlhf,zhou2024archer,wei2025reinforcing} \\
Reward redistribution & holistic scalar re-allocated after the fact & token or time step & \citep{li2024red,qu2024latent} \\
Model-internal attribution & explainability scores or attention flow & token & \citep{koo2025learning,dong2026how} \\
Information-theoretic & belief update about the final answer & turn & \citep{wang2025information} \\
Generative process critique & step-wise critique converted to credit & step, token & \citep{xie2025capo,zhou2024treg,wang2026text2gradreinforcementlearningnatural,liu2025agenticreinforcementlearningimplicit} \\
LLM-as-annotator & subgoal and action judgments by an external LLM & action, subgoal & \citep{pignatelli2024assessing,zhong2024policy,lin2025speaking} \\
Memory-conditioned & credit assigned to memory operations & memory write & \citep{zhou2025memento,yan2026memoryr2} \\
\bottomrule
\end{tabular}
\end{table}

\subsection{Outcome, Process, and Reflective Credit}
\label{subsec:credit-granularity}

Three levels recur. \emph{Outcome credit} scores only the final trajectory. It is cheap and exact for verifiable tasks, which makes it suitable for group-comparison methods, but broadcasting one scalar over every token is maximally coarse \citep{guo2025deepseek}. \emph{Process credit} judges intermediate steps; process supervision and process reward models make a delayed outcome locally actionable \citep{uesato2022solving,lightman2024let}. \emph{Reflective credit} uses the content of an evaluation $\ell$, rather than only its scalarization $\sigma(\ell)$, to determine where and how an update lands. Stored reflections steer later attempts, critiques can be aligned to the spans they indict, and feedback-conditioned policies can retain text as a conditioning variable instead of compressing it \citep{shinn2023reflexion,wang2026text2gradreinforcementlearningnatural,luo2025fcp}. Outcome credit ranks trajectories, process credit ranks steps, and reflective credit can state what to change.

\subsection{Mechanisms and Selection Rules}
\label{subsec:credit-token}

Four mechanisms push sparse outcomes toward the responsible behavior. Temporal abstraction shortens the horizon by treating token spans or turns as macro actions \citep{chai2024marlhf,zhou2024archer}. Reward redistribution preserves a trusted sequence-level scorer but reallocates its value across tokens or time steps \citep{li2024red,qu2024latent}. Model-internal attribution derives token importance from additive explanations or information flow, which makes the shaping signal inspectable but ties it to the reward model's internals \citep{koo2025learning,dong2026how}. Critique-derived methods generate stepwise diagnoses and convert them into token or step credit, gaining semantic localization without a separately labeled dense reward model \citep{xie2025capo,zhou2024treg,liu2025agenticreinforcementlearningimplicit}.

The choice follows the artifacts available before training. Temporal abstraction is the cheapest first move when a trusted outcome reward exists. Redistribution is appropriate when the reward model is the only trusted component; attribution is preferable when the credit trace must be inspected. Critique-derived credit is the only family that can both localize and name an error, but it is justified only when critic quality can be audited. LLM annotation extends the same idea to non-language policies by judging subgoals, transitions, or team contributions, moving the principal failure risk into the annotator \citep{pignatelli2024assessing,zhong2024policy,lin2025speaking}.

\subsection{Long Horizons and Memory}
\label{subsec:credit-multiturn}

Multi-turn interaction introduces two nested questions: which turn helped, and which tokens within it. Hierarchical learners combine turn-level values with within-turn policy gradients; explicit turn rewards are lighter when intermediate outcomes are checkable; information-gain rewards use the policy's belief update when no external reward model is available \citep{zhou2024archer,wei2025reinforcing,wang2025information}. Persistent memory adds a state-conditioned credit problem. Because memory writes change the effective environment seen by later rollouts, group-relative comparisons become biased unless memory operations are compared from a shared intermediate state \citep{zhou2025memento,yan2026memoryr2}. Turn-level rewards are therefore an optional refinement, but memory-conditioned credit is mandatory once the agent writes persistent state.

\subsection{Optimizer Pairing}
\label{subsec:credit-algorithms}

Feedback form constrains the optimizer. Executable rewards pair naturally with online policy gradients; pairwise judgments with offline preference objectives; filtered or feedback-conditioned text with supervised fine-tuning; and step-level judgments with process reward models \citep{schulman2017proximal,rafailov2023direct,luo2025fcp,lightman2024let}. Group comparison adds an assumption that rollouts share an environment and reward scale, which turn structure and memory writes can violate \citep{guo2025deepseek,wei2025reinforcing,yan2026memoryr2}. Language-mediated states also risk information loss, and language-generated actions alter exploration. Planner-imitated subgoals, for example, require explicit exploration to avoid linear regret \citep{he2024words}. The practitioner should therefore choose the feedback carrier first and treat the optimizer as a downstream consequence.

\subsection{The State of Theory}
\label{subsec:credit-theory}

General theory remains thin. The main formal anchor is learning from language feedback through transfer eluder dimension \citep{xu2025formalizing}. Partial results cover regret for language-mediated hierarchies, optimality preservation under feature-additive shaping, redundancy reduction for latent rewards, and the derivation of implicit step rewards \citep{he2024words,koo2025learning,qu2024latent,liu2025agenticreinforcementlearningimplicit}. Missing are conditions under which critique-derived token credit preserves the optimal outcome policy and bounds on how verbal localization errors compound through training. This gap is central to the research agenda in \S\ref{sec:practitioner}.
